\section{Averaging defines a Roe-algebra homomorphism}\label{sec:transfer}

We now prove the analytic input for vanishing. Averaging over each fiber
will embed the module of the filled base discs into the module of $X$.
The main point is to show that the resulting spatial homomorphism takes
values in the locally compact Roe algebra, uniformly over growing fibers.

\subsection{A general transfer criterion}
The following criterion separates that operator argument from the particular
congruence construction. Propagation of an operator between two fiber
subspaces is measured using the ambient metric of $X$.

\begin{proposition}[Transfer through finite fibers]\label{prop:transfer}
Let $D$ and $X$ be countable uniformly discrete bounded-geometry spaces.
Suppose $X=\coprod_{u\in D}F_u$ is a partition into nonempty finite sets.
Write $\xi_u=\abs{F_u}^{-1/2}\sum_{x\in F_u}\delta_x$ and let $P$ be the
orthogonal projection on $\ell^2(X)$ onto the span of the $\xi_u$.
Assume:
\begin{enumerate}
\item $P$ is an operator-norm limit of finite-propagation contractions $P_n$
on $\ell^2(X)$;
\item for every $R<\infty$ there is $s_R<\infty$ such that, whenever
$d_D(u,v)\le R$, there is a contraction
$O_{uv}:\ell^2(F_v)\to\ell^2(F_u)$ of propagation at most $s_R$ with
$O_{uv}\xi_v=\xi_u$.
\end{enumerate}
For the isometry
\[
 V:\ell^2(D)\otimes\HH_0\longrightarrow\ell^2(X)\otimes\HH_0,
 \qquad V(\delta_u\otimes h)=\xi_u\otimes h,
\]
conjugation defines a $*$-homomorphism
$\Ad V:C^*(D)\to C^*(X)$ between the reduced, locally compact Roe algebras.
\end{proposition}
\begin{proof}
The empty case is immediate, so assume $D$ is nonempty. We first lift the
supported blocks of an operator by a finite decomposition into matchings.
Sandwiching this lift between the constant-fiber projections gives the
desired conjugate, and the approximants to those projections place it in
the Roe algebra.

Let $T\in\mathcal B(\ell^2(D)\otimes\HH_0)$ be locally compact and of
propagation at most $R$. Its entries $T_{uv}$ are compact operators on
$\HH_0$. Set $b_R=\sup_u\abs{B_D(u,R)}$. The bipartite support graph of
these entries has degree at most $b_R$, so its edges partition into at most
$N_R=2b_R-1$ matchings. For completeness, enumerate the edges and color
them greedily with $2b_R-1$ colors: at most $2(b_R-1)$ already colored
edges meet the endpoints of the next edge. A matching has no repeated row
or column.

For each matching form the block operator with entries
$O_{uv}\otimes T_{uv}$. Orthogonality of both its source and target blocks
gives norm at most $\sup_{u,v}\norm{T_{uv}}\le\norm T$. Summing the
finitely many matching operators gives a bounded operator $A_T$ with
\begin{equation}\label{eq:matchingnorm}
 \norm{A_T}\le N_R\norm T,\qquad \prop(A_T)\le s_R.
\end{equation}
Its vertex-level entries are scalar multiples of compact operators, so
finite propagation implies local compactness.

Put $\widehat P=P\otimes1_{\HH_0}$ and
$\widehat P_n=P_n\otimes1_{\HH_0}$. Compression of the $(u,v)$ block gives
\[
 p_u O_{uv}p_v=\lvert\xi_u\rangle\langle\xi_v\rvert,
 \qquad \widehat P A_T\widehat P=VTV^*,
\]
where $p_u$ is the rank-one projection onto $\xi_u$. Put
$d_n=\prop(P_n)$. For every $n$,
$\widehat P_n A_T\widehat P_n$ has propagation at most $2d_n+s_R$.
Its $(x,y)$ entry sums only over intermediate vertices
$z\in B_X(x,d_n)$ and $w\in B_X(y,d_n)$. These sets are finite, and
each summand is a scalar multiple of the compact entry $(A_T)_{zw}$.
Thus the sandwich is locally compact and belongs to $C^*(X)$. Moreover,
\begin{equation}\label{eq:sandwich}
 \norm{\widehat P_n A_T\widehat P_n-VTV^*}
 \le 2\norm{P_n-P}N_R\norm T\longrightarrow0.
\end{equation}
It follows that $VTV^*\in C^*(X)$.

Such $T$ are dense in $C^*(D)$, and conjugation by an isometry is norm
continuous. Therefore its range is contained in $C^*(X)$ on the whole
algebra. Finally $V^*V=1$ proves multiplicativity and preservation of
adjoints. The auxiliary choices in $A_T$ are used only to prove this range
inclusion; they need not depend linearly on $T$.
\end{proof}

The amplified projection $\widehat P$ in this proof is a bounded operator,
not a locally compact Roe-algebra element. Its diagonal compression at
$x\in F_u$ is $\abs{F_u}^{-1}1_{\HH_0}$, which is not compact.
The proof establishes local compactness for the sandwiches, where it is
actually needed.

\subsection{Uniform approximation from property (T)}
Return to the finite fibers $F_u=\{u\}\times Q_{k_j(u)}$ of our graph $X$.
Kazhdan introduced property~(T) and proved its higher-rank and lattice
instances \cite{Kazhdan}. For $\SL_3(\Z)$ it provides
a constant $\kappa>0$ for the fixed
generating set $S$: in any unitary representation without invariant vectors,
each unit vector is moved by at least $\kappa$ by some $s\in S$
\cite[Proposition~1.3.2 and Example~1.7.4]{BHV}.

Let $\lambda_k$ denote the left regular representation on $\ell^2(Q_k)$,
and define
\[
 L_k=1-\frac1{\abs S}\sum_{s\in S}\lambda_k(s).
\]
The sum retains multiplicities when elements of $S$ have the same image
in $Q_k$. Since $S$ is symmetric, $0\le L_k\le2$. Its kernel consists
exactly of the constants, since $S$ generates $Q_k$. On the orthogonal
complement of the constants,
\[
 \langle L_k\eta,\eta\rangle
 =\frac1{2\abs S}\sum_{s\in S}\norm{\lambda_k(s)\eta-\eta}^2
 \ge\delta\norm\eta^2,\qquad \delta=\frac{\kappa^2}{2\abs S}>0.
\]
If $p_k$ denotes projection onto the normalized constant vector, it follows
that
\begin{equation}\label{eq:gap}
 \norm{(1-L_k/2)^n-p_k}\le(1-\delta/2)^n.
\end{equation}
For $k=0$ the approximation is exact. Take the direct sum of these lazy
averaging powers over all fibers and call it $P_n$. They are contractions
of propagation at most $n$, since each Cayley generator moves at most one
edge in $X$. Equation~\eqref{eq:gap} proves uniform norm convergence to
the constant-fiber projection $P$. This verifies the first hypothesis of
Proposition~\ref{prop:transfer}.

\subsection{Controlled operators between fibers}
It remains to verify the second hypothesis. Equal-depth edges act by the
permutation prescribed by their right holonomy. For an edge from a shallow
vertex $u$ of depth $k$ to a deep vertex $v$ of depth $k+1$, put
$a=a(u,v)$ and $h_k=\abs{\ker(Q_{k+1}\to Q_k)}$. The graph relation gives
an isometry $J_a:\ell^2(Q_k)\to\ell^2(Q_{k+1})$ and its adjoint:
\begin{equation}\label{eq:childmap}
 J_a\delta_q=h_k^{-1/2}
       \sum_{\bar z=qg_k^a}\delta_z,
 \qquad
 J_a^*\delta_z=h_k^{-1/2}\delta_{\bar z g_k^{-a}}.
\end{equation}
Every coefficient lies on a graph edge, so these maps have propagation at
most one. The child sets have the common cardinality $h_k$ and are disjoint;
this proves the isometry assertion. Since $\abs{Q_{k+1}}=h_k\abs{Q_k}$,
both $J_a$ and its adjoint carry the normalized constant vectors to the
corresponding normalized constants: each deep coefficient in the image
under $J_a$ is $(h_k\abs{Q_k})^{-1/2}$, and each shallow coefficient
in the image under $J_a^*$ is $h_k(h_k\abs{Q_{k+1}})^{-1/2}$.
These formulas also apply to the singleton quotient $Q_0$.
The equal-depth permutations have the
same property.

For $u,v$ in the same base disc with $d_D(u,v)\le R$, choose a base path
from $v$ to $u$ of length at most $\lfloor R\rfloor$. Composing the edge
operators just described gives a contraction $O_{uv}$ of propagation at
most $\lfloor R\rfloor$, with $O_{uv}\xi_v=\xi_u$. This includes $u=v$
by using the identity on its fiber.

For the cross-component pairs at distance at most $R$, use
$O_{uv}=\lvert\xi_u\rangle\langle\xi_v\rvert$. There are only finitely
many such pairs by Section~\ref{sec:geometry}, and their fibers are finite.
Consequently their propagations have a common finite bound. Combining this
bound with $\lfloor R\rfloor$ supplies $s_R$ in
Proposition~\ref{prop:transfer}. We conclude that
\begin{equation}\label{eq:Phi}
 \Phi=\Ad V:C^*(D)\longrightarrow C^*(X)
\end{equation}
is a $*$-homomorphism.
This conclusion is an operator-norm statement in the reduced Roe algebra.
It does not require $V$ to cover a coarse map from $D$ to $X$.
